В третьей главе рассматривается практическая реализация вычислительного модуля и его интеграция с вспомогательными компонентами проекта. Основной акцент сделан на устройстве пакета \texttt{compute}, поскольку именно он реализует математический аппарат адаптивных разреженных сеток, поддержку линейного и квадратичного базисов, параллельное построение, контроль числа узлов и механизмы работы с динамическими системами.

\subsection{Реализация вычислительного ядра на Go}

Основная часть алгоритмов сосредоточена в пакете \texttt{compute}. Его центральным элементом является тип \texttt{AdaptiveSparseGrid}, инкапсулирующий параметры области, тип базиса, размерности входа и выхода, список входных точек и множество узлов. Выбор Go в качестве языка реализации вычислительного модуля обусловлен поддержкой конкурентных вычислений, удобной работой с функциями высшего порядка и простотой интеграции в прикладной сервис \cite{GoDocs}.

Конструктор сетки выполняет валидацию входных данных, определяет размерность образа функции, вычисляет значения в контрольных точках и затем запускает построение. Важной особенностью является то, что вычислительное ядро принимает пользовательскую функцию как аргумент. Благодаря этому один и тот же код используется и для интерполяции аналитических отображений, и для операторов перехода, получаемых численным интегрированием.

В листинге \ref{lst:new_grid} показан фрагмент конструктора, задающего общий сценарий создания сетки.

\begin{lstlisting}[language=Go, caption={Создание адаптивной разреженной сетки}, label={lst:new_grid}]
func NewAdaptiveSparseGridWithContext(
    ctx context.Context,
    funcEval func(Point) (Point, error),
    min, max Point,
    epsilon ScalarType,
    anchorPoints []Point,
    basisType BasisType,
    buildType BuildType,
    maxLevel int64,
    maxNodesInGrid int64,
) (*AdaptiveSparseGrid, error) {
    inDim := int64(len(min))
    if int64(len(max)) != inDim {
        return nil, fmt.Errorf("min and max must have the same dimension")
    }

    sampleOut, err := funcEval(min)
    if err != nil {
        return nil, err
    }
    outDim := int64(len(sampleOut))

    grid := &AdaptiveSparseGrid{
        min:       min,
        max:       max,
        basisType: basisType,
        inDim:     inDim,
        outDim:    outDim,
        nodes:     make(map[string]node),
    }

    anchors, err := grid.calculateOriginalFunctionAtPoints(funcEval, anchorPoints)
    if err != nil {
        return nil, err
    }

    err = grid.build(ctx, funcEval, epsilon, anchors, buildType, maxLevel, maxNodesInGrid)
    if err != nil {
        return nil, err
    }

    return grid, nil
}
\end{lstlisting}

В данном фрагменте видны основные свойства реализации: независимость от природы исходной функции, автоматическое определение размерности результата и наличие параметров, отвечающих за точность, тип базиса, режим построения и жёсткое ограничение числа узлов.

\subsection{Ключи узлов и кодирование положения в сетке}

Каждый узел сетки идентифицируется набором уровней и индексов. Для хранения узлов в отображении используется тип \texttt{gridKey}. Его строковое представление формируется бинарной упаковкой целочисленных массивов.

В листинге \ref{lst:grid_key} показан соответствующий код.

\begin{lstlisting}[language=Go, caption={Представление ключа узла сетки}, label={lst:grid_key}]
type gridKey struct {
    level []int64
    index []int64
}

func (k gridKey) String() string {
    totalLength := (len(k.level) + len(k.index)) * 8
    result := make([]byte, totalLength)

    pos := 0
    for _, v := range k.level {
        binary.LittleEndian.PutUint64(result[pos:], uint64(v))
        pos += 8
    }
    for _, v := range k.index {
        binary.LittleEndian.PutUint64(result[pos:], uint64(v))
        pos += 8
    }

    return unsafe.String(&result[0], len(result))
}
\end{lstlisting}

Достоинство этого решения состоит в быстром поиске узла по его иерархическому положению. Это критически важно для построения потомков и для рекурсивной интерполяции по дереву узлов.

\subsection{Реализация базисных функций}

Математическая модель из первой главы непосредственно отражается в коде вычисления одномерного базиса. В зависимости от уровня, индекса и выбранного типа базиса вычисляется значение локальной функции влияния. Именно здесь программно различаются линейный и квадратичный варианты.

В листинге \ref{lst:eval1d} приведена функция вычисления одномерного базиса.

\begin{lstlisting}[language=Go, caption={Вычисление одномерной базисной функции}, label={lst:eval1d}]
func eval1d(x float64, level int64, index int64, basisType BasisType) float64 {
    if level == 0 {
        if index == 0 {
            return 1.0 - x
        }
        if index == 1 {
            return x
        }
        return 0.0
    }

    h := 1.0 / float64(int64(1)<<level)
    center := float64(index) * h
    dist := math.Abs(x - center)

    if dist >= h {
        return 0.0
    }

    t := dist / h
    switch basisType {
    case BasisTypeLinear:
        return 1.0 - t
    case BasisTypeQuadratic:
        return 1.0 - t*t
    }
    return 0.0
}
\end{lstlisting}

Данный фрагмент важен не только как реализация базиса, но и как точка прямого сопоставления двух режимов работы, сравнение которых затем проводится в бенчмарках.

\subsection{Параллельный и последовательный режимы построения}

Поддержка двух режимов сборки реализована без раздвоения математической логики. Независимые стартовые конфигурации либо вычисляются по очереди, либо запускаются как отдельные goroutine. На практике это позволяет сравнивать режимы \texttt{parallel} и \texttt{sequential} на одном и том же наборе задач.

В листинге \ref{lst:build_parallel} приведён фрагмент функции \texttt{build}, показывающий выбор между последовательным и параллельным режимами.

\begin{lstlisting}[language=Go, caption={Параллельный запуск локальных задач построения}, label={lst:build_parallel}]
if buildType == BuildTypeParallel {
    wg.Add(1)
    go func(k gridKey) {
        defer wg.Done()
        res := g.buildGridTask(ctx, funcEval, epsilon, anchors, k, i, maxLevel, maxNodesInGrid)
        mu.Lock()
        results = append(results, res)
        mu.Unlock()
    }(currentKey)
} else {
    res := g.buildGridTask(ctx, funcEval, epsilon, anchors, currentKey, i, maxLevel, maxNodesInGrid)
    results = append(results, res)
}
\end{lstlisting}

Такой способ организации вычислений хорошо согласуется с логикой метода, поскольку параллелизуются именно независимые задачи первого уровня.

\subsection{Создание узла и вычисление коэффициента избытка}

При создании нового узла вычисляется значение исходной функции, затем значение текущей интерполяции в той же точке, после чего формируется коэффициент избытка как их разность.

В листинге \ref{lst:create_node} показан соответствующий фрагмент кода.

\begin{lstlisting}[language=Go, caption={Формирование узла и вычисление коэффициента избытка}, label={lst:create_node}]
func (g *AdaptiveSparseGrid) createNode(
    funcEval func(Point) (Point, error),
    key gridKey,
    entryPoint *node,
    dimension int64,
    additionalNodes map[string]node,
) (*gridKey, error) {
    if _, exists := additionalNodes[key.String()]; exists {
        return nil, nil
    }

    n := node{
        key:        key,
        centerUnit: make(Point, g.inDim),
    }
    for i := int64(0); i < g.inDim; i++ {
        n.centerUnit[i] = getCoord(key.level[i], key.index[i])
    }

    realArg := g.toReal(n.centerUnit)
    etalon, err := funcEval(realArg)
    if err != nil {
        return nil, err
    }

    interp := g.evaluateForDimAndEntryPoint(n.centerUnit, dimension, entryPoint, additionalNodes)
    etalon.Sub(interp)
    n.alpha = etalon
    additionalNodes[key.String()] = n

    return &key, nil
}
\end{lstlisting}

Этот код непосредственно реализует иерархическую природу метода: узел хранит не абсолютное значение функции, а именно добавочную информацию относительно уже построенного приближения.

\subsection{Ключевой участок buildGrid: адаптация, anchor points и ограничение числа узлов}

Наиболее важным фрагментом реализации является внутренняя логика \texttt{buildGrid}, где одновременно решаются задачи анализа коэффициента избытка, учёта anchor points, порождения дочерних узлов и контроля числа узлов. Именно этот участок лучше всего отражает реальное содержание вычислительного модуля, поэтому он является более показательным, чем элементарные операции над типом \texttt{Point}.

В листинге \ref{lst:build_grid_core} приведён характерный фрагмент алгоритма.

\begin{lstlisting}[language=Go, caption={Фрагмент алгоритма buildGrid: критерии продолжения и ограничение числа узлов}, label={lst:build_grid_core}]
for len(nodeQueue) > 0 {
    currentKey := nodeQueue[0]
    nodeQueue = nodeQueue[1:]

    currentNode := newNodes[currentKey.String()]
    activeEntryNode := newNodes[entryPointKeyString]

    canContinue := false
    directions := make([]direction, g.inDim)
    for i := range directions {
        directions[i] = directionBoth
    }

    if currentNode.alpha.Length() <= epsilon {
        for i := range directions {
            directions[i] = directionNone
        }
        canContinue = true
    }

    if canContinue {
        for _, anchor := range anchors {
            if currentNode.isPointInAffectZone(anchor.arg, g.inDim) {
                evalRes := g.evaluateForDimAndEntryPoint(anchor.arg, dimension, &activeEntryNode, newNodes)
                evalRes.Sub(anchor.ans)
                if evalRes.Length() >= epsilon {
                    canContinue = false
                    affectDirs := getAffectDirection(anchor.arg, currentNode.key.level, currentNode.key.index, g.inDim)
                    for i := int64(0); i < g.inDim; i++ {
                        directions[i] |= affectDirs[i]
                    }
                }
            }
        }
    }

    if canContinue {
        continue
    }

    currentNode.hasChildren = true
    newNodes[currentKey.String()] = currentNode

    if maxNodesInGrid > 0 && int64(len(newNodes)) >= maxNodesInGrid {
        break
    }
}
\end{lstlisting}

Этот фрагмент отражает несколько существенных особенностей реализации:
\begin{enumerate}
    \item базовое решение принимается по норме коэффициента избытка;
    \item anchor points могут отменить преждевременную остановку;
    \item уточнение выполняется направленно;
    \item рост сетки может быть жёстко ограничен параметром \texttt{maxNodesInGrid}.
\end{enumerate}

Именно здесь проявляется одно из главных отличий реализованного модуля от базовой математической схемы, рассматриваемой в статье \cite{Morozov2024}: введение ручного механизма управления уточнением через anchor points.

\subsection{Рекурсивная интерполяция по построенной сетке}

После завершения построения объект \texttt{AdaptiveSparseGrid} используется для быстрого вычисления приближённых значений. Интерполяция организована как рекурсивный обход дерева узлов: сначала учитывается вклад текущего узла, затем --- вклад релевантных потомков.

В листинге \ref{lst:evaluate_recursive} показан характерный фрагмент рекурсивного вычисления.

\begin{lstlisting}[language=Go, caption={Рекурсивное вычисление интерполянта}, label={lst:evaluate_recursive}]
func (g *AdaptiveSparseGrid) evaluateRecursive(
    x Point,
    n node,
    processedNodes map[string]struct{},
    additionalNodes map[string]node,
) Point {
    if _, exists := processedNodes[n.key.String()]; exists {
        return make(Point, g.outDim)
    }

    processedNodes[n.key.String()] = struct{}{}
    basis := evalBasis(x, n.key.level, n.key.index, g.basisType, g.inDim)

    if math.Abs(basis) < 1e-9 {
        return make(Point, g.outDim)
    }

    value := n.alpha.Clone()
    value.MulScalar(basis)

    if n.hasChildren {
        for i := int64(0); i < g.inDim; i++ {
            if n.key.level[i] != 0 {
                child, exists := g.getChildForDimAndArg(n, x, i, additionalNodes)
                if exists {
                    value.Add(g.evaluateRecursive(x, child, processedNodes, additionalNodes))
                }
            }
        }
    }
    return value
}
\end{lstlisting}

Программный смысл этого кода совпадает с теоретической формулой суммы иерархических вкладов. Если значение базиса практически нулевое, соответствующая ветвь обхода мгновенно отсекается.

\subsection{Два режима работы с дифференциальными уравнениями}

Для задач динамики в проекте поддерживаются два принципиально разных подхода.

\textbf{Первый подход: параметр $t$ как интервальная неопределённость.} В этом случае время добавляется как ещё одна координата пространства аргументов. Тогда строится сетка для отображения
\begin{equation}
    (x_0, t) \mapsto x(t).
\end{equation}

Преимущество этого подхода состоит в том, что одна построенная сетка сразу описывает поведение решения на целом временном интервале. Недостаток --- рост размерности задачи и, как следствие, рост вычислительной сложности.

\textbf{Второй подход: итеративное построение через \texttt{MakeNextIteration}.} В этом случае сначала строится одношаговый оператор перехода на интервале длины $1$, а затем он применяется последовательно несколько раз. Преимущество такого подхода состоит в том, что размерность пространства не увеличивается, но приходится выполнять серию последовательных построений.

Именно сопоставление этих двух подходов является одним из ключевых содержательных результатов экспериментальной части, поскольку оно показывает, что разные трактовки времени приводят к разным профилям вычислительной сложности.

\subsection{Вспомогательный серверный слой}

Пакет \texttt{internal/server} реализует удалённый интерфейс к вычислительному ядру. Объект \texttt{Formula} использует библиотеку \texttt{expr} для интерпретации пользовательских выражений \cite{ExprLang}. При работе с дифференциальными уравнениями сервер численно интегрирует систему методом Рунге--Кутты--Фельберга \cite{Hairer1993}, после чего передаёт в модуль \texttt{compute} уже готовое отображение.

С точки зрения архитектуры серверный слой является адаптером, а не предметом основной новизны работы. Его значение состоит в том, что он позволяет использовать вычислительный модуль удалённо и удобно подключать визуализацию.

\subsection{Клиентская визуализация}

Клиентская часть на Flutter реализует ввод параметров, запуск вычислений и отображение узлов сетки \cite{FlutterDocs}. Она важна как средство анализа результатов, но не является математическим ядром работы. Основной результат диплома по-прежнему сосредоточен в модуле \texttt{compute}; клиент выступает как инструмент демонстрации и визуального исследования структуры построенного интерполянта.

\subsection{Анализ структуры бенчмарков}

Для оценки эффективности реализации был расширен набор бенчмарков. В текущей версии проекта измеряются:
\begin{enumerate}
    \item зависимость времени построения от размерности задачи и размера области;
    \item сравнение последовательного и параллельного режимов построения;
    \item сравнение линейного и квадратичного базисов при различных значениях $\varepsilon$;
    \item влияние anchor points на сложность построения;
    \item сравнение двух подходов к работе с дифференциальными уравнениями;
    \item влияние жёсткого ограничения числа узлов;
    \item стоимость вычисления значения по уже построенной сетке.
\end{enumerate}

Такой набор экспериментов нельзя считать тривиальным, поскольку он отражает не только абсолютное время работы, но и влияние архитектурных решений модуля \texttt{compute} на разные режимы использования.

\subsection{Анализ полученных данных}

По содержанию бенчмарков можно выделить несколько наиболее важных групп выводов.

Во-первых, линейный и квадратичный базисы должны сравниваться не только по качеству аппроксимации, но и по времени построения и стоимости последующего вычисления. Квадратичный базис обычно оказывается более затратным на локальном уровне, но при этом позволяет получать более гладкое приближение и потенциально уменьшать необходимую глубину уточнения. Поэтому содержательный анализ должен учитывать обе стороны --- локальную стоимость и глобальную эффективность.

Во-вторых, сравнение \texttt{parallel} и \texttt{sequential} показывает практическую ценность распараллеливания стартовых конфигураций. Это сравнение особенно важно для темы диплома, поскольку напрямую связано с заявленной реализацией параллельных алгоритмов.

В-третьих, сравнение двух режимов работы с дифференциальными уравнениями демонстрирует, что параметр времени может быть учтён либо как отдельная координата, либо через последовательную композицию одношагового оператора. Эти подходы различаются по архитектурному смыслу и по производительности, а потому требуют отдельного разбора в итоговой аналитической главе.

В-четвёртых, механизм жёсткого ограничения числа узлов оказывается важным компромиссом между точностью и затратами. Он даёт возможность использовать метод в ресурсно ограниченной среде и должен рассматриваться как существенный элемент практической оптимизации.

\subsection*{Выводы по главе}
\addcontentsline{toc}{subsection}{Выводы по главе}

В третьей главе была рассмотрена практическая реализация вычислительного модуля и его вспомогательной инфраструктуры. Показано, каким образом математический аппарат разреженных сеток отражён в коде пакета \texttt{compute}, как реализованы базисные функции, коэффициенты избытка, адаптивное уточнение, anchor points, ограничение числа узлов, рекурсивная интерполяция и два режима работы с дифференциальными уравнениями. Отдельно рассмотрены архитектурная роль серверного слоя и клиентской визуализации, а также структура расширенного набора бенчмарков. Всё это подтверждает, что основное содержание дипломной работы сосредоточено именно в вычислительном модуле, а дополнительные компоненты служат средствами его практического применения и анализа.
